% appE —— 附录 E 与热力学贝特拟设的关系（原文 wcf-paper.tex 3900–4041 行）
\section{与热力学贝特拟设的关系}
\label{app:TBA}

{\itshape 2009 年 11 月 20 日补注}：

不久前，A. Zamolodchikov 向我们指出：本文的核心结果之一，即式 (5.13)，实际上是热力学贝特拟设（Thermodynamic Bethe Ansatz，下称 TBA）的一个版本 \cite{Zamolodchikov:1989cf}。本附录解释这一评注。四维超杨--米尔斯（super Yang--Mills）理论与 TBA 之间的另一个联系，最近也已由 Nekrasov 与 Shatashvili 讨论过 \cite{Nekrasov:2009ui}。% TODO 合稿改 \eqref{eq:X-integral-mult-explicit}

考虑一个可积的粒子系统：粒子记作 $a$，质量为 $m_a$，处于逆温度 $\beta$ 下，散射矩阵 $S_{ab}(\theta-\theta')$ 可积，其中 $\theta$ 为快度（rapidity）。该系统的 TBA 方程为
%
\begin{equation}\label{eq:gen-TBA}
\epsilon_a(\theta) = m_a \beta \cosh \theta - \sum_b
\int_{-\infty}^{+\infty} \frac{d \theta'}{2\pi}
\phi_{ab}(\theta-\theta') \log (1 + e^{\beta \mu_b -
\epsilon_b(\theta')})
\end{equation}
%
其中 $\phi_{ab}(\theta) = - i\frac{\p}{\p \theta} \log
S_{ab}(\theta)$。这里散射矩阵是对角的，也就是说，孤子产生算子满足 $\Phi_a(\theta) \Phi_b(\theta') =
S_{ab}(\theta-\theta') \Phi_b(\theta') \Phi_a(\theta)$。

我们可以把式 (5.13) 的对数写成 \eqref{eq:gen-TBA} 的形式，做法如下。显然，粒子标号 $a, b, \dots$ 对应于 $\gamma, \gamma', \dots$。现在令% TODO 合稿改 \eqref{eq:X-integral-mult-explicit}（式 (5.13)）
$Z_\gamma =
e^{i \alpha_\gamma} \vert Z_\gamma \vert$，其中 $\alpha_\gamma$ 为实数，且只在模 $2\pi$ 的意义下有定义。对任意一个 $\gamma$，我们都可以作变量替换 $\zeta = - e^{i \alpha_\gamma+ \theta}$，使 BPS 射线 $\ell_\gamma$ 由 $-\infty < \theta <
\infty$ 扫出。在这一变量替换下，半平坦坐标式 (3.9) 变为% TODO 合稿改 \eqref{eq:X-integral-mult-explicit}（式 (5.13)）、\eqref{eq:X-sf}（式 (3.9)）
%
\begin{equation}
\log \CX_{\gamma}^{\sf} = - 2 \pi R \lvert Z_\gamma \rvert \cosh \theta
+ i \varphi_\gamma.
\end{equation}
%
（注意：为了避免与快度 $\theta$ 混淆，我们把环面上角坐标的记号，从本文其余部分所用的 $\theta_\gamma$ 改成了 $\varphi_\gamma$。）现在置
%
\begin{equation}\label{eq:defmugamma}
\beta \mu_\gamma := i \varphi_\gamma + \log(-\sigma(\gamma))  \quad \mod \ 2
\pi i.
\end{equation}
%
注意，$\beta \mu_\gamma$ 就是 $i \varphi_\gamma$，或者与之相差 $\pm i \pi$。特别地，它是纯虚数。我们按下式定义「准粒子（quasiparticle）能量」$\epsilon_{\gamma}(\theta)$：
%
\begin{equation} \label{eq:qparticle}
\CX_{\gamma}(\zeta = - e^{i \alpha_\gamma+ \theta}) :=
-\sigma(\gamma) e^{\beta \mu_\gamma - \epsilon_{\gamma}(\theta)}=
e^{i \varphi_\gamma - \epsilon_\gamma(\theta)}.
\end{equation}
%
更确切地说，上式所定义的是 BPS 射线 $\ell_\gamma$ 上的 $\eps_\gamma$，在那里 $\theta$ 取实值。对其他 $\theta$，我们用解析延拓来定义 $\eps_\gamma$——这与 $\CX_\gamma$ 恰成对照：后者沿着某些 $\Im\,\theta$ 为常数的直线（即 BPS 射线）带有不连续性。

我们选取 \eqref{eq:defmugamma}，正是为了使式 (5.13) 的对数读作% TODO 合稿改 \eqref{eq:X-integral-mult-explicit}（式 (5.13)）
%
\begin{equation}\label{eq:GMN-2}
\epsilon_{\gamma}(\theta) = 2\pi R \vert Z_\gamma \vert \cosh\theta
+ \sum_{\gamma'} \Omega(\gamma') \int_{-\infty}^{+\infty}
\frac{d\theta'}{2\pi} K_{\gamma,\gamma'}(\theta-\theta') \log( 1 +
e^{\beta \mu_{\gamma'} - \epsilon_{\gamma'}(\theta')} ),
\end{equation}
%
其中
%
\begin{equation}
K_{\gamma,\gamma'}(\theta-\theta') = \frac{i}{2} \langle \gamma,
\gamma'\rangle \frac{ e^{ \theta - \theta' + i \alpha_{\gamma} - i
\alpha_{\gamma'}} +1 }{e^{ \theta - \theta' + i \alpha_{\gamma} - i
\alpha_{\gamma'}} -1 }.
\end{equation}
%
这个核也可以写成
%
\begin{equation}
\begin{aligned}
K_{\gamma,\gamma'}(\theta-\theta') & = i \langle \gamma, \gamma'
\rangle \frac{\p}{\p \theta} \log \bigl[\sinh\left(\half (\theta -
\theta' + i \alpha_{\gamma} - i \alpha_{\gamma'})\right)\bigr]\\
& =\frac{i}{2}\langle \gamma, \gamma' \rangle  \coth\left( \frac{
\theta-\theta' + i \alpha_{\gamma} - i \alpha_{\gamma'}}{2} \right).
\end{aligned}
\end{equation}
%
对数的宗量并不是一个纯相位，因此一般说来，$K_{\gamma,\gamma'}(\theta-\theta')$ 并不对应于一个幺正的散射矩阵。

下面简略讨论「准粒子能量」的实性（reality）性质。扭量坐标满足如下实性关系
%
\begin{equation}
\overline{\CX_{\gamma}(\zeta)} = \CX_{-\gamma}(-1/\bar \zeta).
\end{equation}
%
由于 $Z_{-\gamma} = - Z_\gamma$，我们有
%
\begin{equation}
e^{i \alpha_{-\gamma}} = - e^{i \alpha_\gamma}.
\end{equation}
%
因此，若 $\zeta = - e^{i \alpha_{\gamma} + \theta}$ 且 $\theta$ 为实数，则 $-1/\bar \zeta = - e^{i \alpha_{-\gamma} - \theta}$。现在，利用 $\sigma(-\gamma) = \sigma(\gamma)$ 与 $\varphi_{-\gamma} = -
\varphi_{\gamma}$，便得到「准粒子能量」的实性条件
%
\begin{equation}\label{eq:reality}
\overline{\epsilon_{\gamma}(\theta)} = \epsilon_{-\gamma}(-\theta).
\end{equation}
%
积分方程式 (5.13) 与实性条件 \eqref{eq:reality} 相容，因为当 $\theta, \theta'$ 均为实数时% TODO 合稿改 \eqref{eq:X-integral-mult-explicit}（式 (5.13)）
%
\begin{equation}
\overline{K_{\gamma, \gamma'}(\theta-\theta')} = K_{- \gamma, -
\gamma'}(-\theta + \theta').
\end{equation}

\begin{yizhu}
\begin{enumerate}[leftmargin=2em,itemsep=2pt,topsep=2pt]
\item \textbf{关于 $\sigma(\gamma)$}：\eqref{eq:defmugamma} 中的 $\sigma(\gamma)$ 是第 2.2 节引入的「二次加细」，即满足 $\sigma(\gamma_1)\sigma(\gamma_2) = (-1)^{\langle \gamma_1, \gamma_2\rangle} \sigma(\gamma_1+\gamma_2)$ 的函数 $\sigma:\clat\to\IZ_2$，故取值 $\pm1$；于是 $\log(-\sigma(\gamma))$ 只能取 $0$ 或 $\pm i\pi$，这正是原文「$\beta\mu_\gamma$ 就是 $i\varphi_\gamma$，或者与之相差 $\pm i\pi$」那句话的来由。% TODO 合稿改 \ref{sec:KS-formula}（第 2.2 节）
\item \textbf{与 \eqref{eq:gen-TBA} 的对照}：把 \eqref{eq:gen-TBA} 与 \eqref{eq:GMN-2} 并排比较时应注意，两式的约定并不逐字相同：\eqref{eq:gen-TBA} 对粒子种类 $b$ 求和，且 $\phi$ 之前带负号；\eqref{eq:GMN-2} 则对电荷 $\gamma'$ 求和、以 BPS 简并度 $\Omega(\gamma')$ 为权重，且 $K$ 之前带正号。
\end{enumerate}
\end{yizhu}
